\documentclass[12pt]{article}

\usepackage{amssymb, color}
\usepackage[cmex10]{amsmath}
\usepackage{xcolor,graphicx}
\usepackage[margin=1in]{geometry}
\usepackage{fancyhdr}
\usepackage{enumitem}

\usepackage[pdftex,bookmarks,letterpaper,hyperfigures,colorlinks
						,urlcolor=blue
						,citecolor=blue
						,linkcolor=blue
						,pagecolor=blue
						,pdfstartview=FitH]{hyperref}
\providecommand{\hreff}[1]{\href{http://#1}{http://#1}}
\providecommand{\email}[1]{\href{mailto:#1}{\textcolor{black}{#1}}}
\newcommand{\bbm}{\begin{bmatrix}}
\newcommand{\ebm}{\end{bmatrix}}

% Setup style for fancyheader
\pagestyle{fancyplain} % or fancyplain, plain

\lfoot{\textit{Last Modified: \today}}
\cfoot{}
\rfoot{\thepage}

\begin{document}
\begin{center}{\Large \bf{APMTH 50 Problem Set 5}} \vspace{0.2in}\\
{\large Due on Friday, April 21, by 5pm\\}
\end{center}

\subsection*{Collaboration policy}
Collaboration on homework and labs is
encouraged, but each student must write solutions in their own words, and not
copy them from anyone else.  Please write the name(s) of your collaborator(s) at the top of the homework.

\subsection*{Submission guidelines} Please upload your homework as a single Word or pdf document.   \textbf{For this homework, you should submit all of your code as well as the output from your code.}%and but you should include all output (such as figures/plots) from your code.  
\subsection*{Final project workshops} 
Classes this week will be final project workshops.  Please come prepared to discuss your project and what you have done on it so far.  See comments on your homework 4 Canvas submission for ideas on how to  prepare for these workshops.

\subsection*{Problems}
\begin{enumerate}
 \item \textbf{Unique number strategy.} Submit your Matlab function for the unique
    number competition, following the instructions and information from the
    lab. Your routine should be called \texttt{firstname\_lastname.m} and
    should be uploaded as a separate file to Canvas.
  \item \textbf{Rock-paper-scissors.} In this
    \href{http://en.wikipedia.org/wiki/Rock-paper-scissors}{well-known game},
    two players simultaneously choose between the options of rock, paper, or
    scissors. The winner is then decided by using the following conditions:
    rock beats scissors, scissors beats paper, and paper beats rock.
    \begin{enumerate}
      \item Suppose all that matters is the total number of wins, so that the
	payoff matrix for the first player is
	\begin{equation}
	  A=\left(
	  \begin{array}{ccc}
	    0 & 0 & 1 \\
	    1 & 0 & 0 \\
	    0 & 1 & 0
	  \end{array}	  
	  \right).
	  \label{eq:rps1}
	\end{equation}
	Consider playing a mixed strategy where rock, paper, and scissors are
	chosen with probabilities $p_0$, $p_1$, and $p_2$ respectively. Calculate
	the expected payoff for playing either rock, paper, or scissors. Use
	these to determine a set of three differential equations for $p_0$, $p_1$,
	and $p_2$ for replicator dynamics following the procedure described in
	the lab. Starting from $(p_0,p_1,p_2)=(0.8,0.1,0.1)$, solve the system
	of differential equations numerically to find equilibrium values of
	$p_0$, $p_1$, and $p_2$.
      \item Now suppose that a win scores one point, and a loss loses one
	point, so that the payoff matrix for the first player is
	\begin{equation}
	  A=\left(
	  \begin{array}{ccc}
	    0 & -1 & 1 \\
	    1 & 0 & -1 \\
	    -1 & 1 & 0
	  \end{array}
	  \right).
	  \label{eq:rps2}
	\end{equation}
	Modify your program from part (a) to solve the replicator dynamics for
	this payoff matrix. Try the following two initial conditions: (i) your
	equilibrium solution from part (a), and (ii)
	$(p_0,p_1,p_2)=(0.8,0.1,0.1)$. Comment on the results.
      \item \textbf{Extra credit.} Examine the case where a win scores one point and
	a loss loses two points.
    \end{enumerate}
%  \item \textbf{Rock-paper-scissors against a general opponent.} Suppose that you
%    play rock-paper-scissors against an opponent who chooses the three options
%    with probability $q_0$, $q_1$, and $q_2$. Assume that these probabilities
%    may not represent an optimal strategy, and could take any values, provided
%    that $q_0+q_1+q_2=1$.
%    \begin{enumerate}
%      \item Suppose that you play a mixed strategy where you choose the three
%	options with probabilities $p_0$, $p_1$, and $p_2$ where
%	$p_0+p_1+p_2=1$. Write down the expected payoff $P_a(p_0,p_1,p_2)$
%	corresponding to the payoff matrix given in question 2(a). What values
%	of $p_0$, $p_1$, and $p_2$ will maximize $P_a$?
%      \item Suppose that you are interested in maximizing the payoff $P_b$
%	corresponding to the matrix in question 2(b). What values of $p_0$,
%	$p_1$, and $p_2$ will maximize $P_b$?
%      \item Construct an example set of values for $q_0$, $q_1$, and $q_2$
%	where the optimal strategies in parts (a) and (b) will differ.
%    \end{enumerate}
  \item \textbf{Replication dynamics for a small unique number game.}
    Consider a unique number game where three players A, B, and C choose from
    the numbers 0, 1, and 2.
    \begin{enumerate}
      \item Suppose that A chooses 0. What are the possible choices of the
	other two players that will lead to A winning? Repeat this for
	when A chooses 1, and when A chooses 2.
      \item Suppose that the three players are employing a mixed strategy in
	which the numbers 0, 1, and 2 are chosen with probabilities $p$, $q$, and
	$r$ respectively. Suppose that winning a round scores one, and that all
	other outcomes score zero. By making use of your results from part (a),
	calculate the expected payoff of choosing 0, 1, or 2.
      \item Use your results from part (b) to set up a replicator dynamics
	system of three differential equations for $p$, $q$, and $r$ following
	the procedure described in the lab. Solve the system and determine
	an equilibrium solution, making sure that you simulate the system over
	a long enough time interval to obtain three digits of precision in the
	equilibrium probabilities.
      \item \textbf{Extra credit.} Repeat the analysis in steps (a), (b), and (c) for the following extensions to the above scenario:
	\begin{enumerate}
	  \item Four players choosing from 0, 1, and 2.
	  \item Three players choosing from 0, 1, 2, and 3.
	  \item Three players choosing from 0, 1, and 2, where a player scores
	   $ \frac{1}{3}$ when no unique number is chosen.
	\end{enumerate}
      \item \textbf{Extra credit.} Suppose that you are playing many repeated games,
	and your objective is to win more games than your opponents. What is a
	better strategy: your answer from (c) or your answer from (d)(iii)?
    \end{enumerate}

\end{enumerate}
\end{document}


